\section{Architecture and portability}
\begin{tabular}{p{0.30\linewidth}p{0.62\linewidth}}
\toprule Location & Responsibility \\ \midrule
Core & CubeMX-owned main, interrupts, HAL configuration and system clock. \\
Application & Board adapter, DWT delay, GPIO callbacks and example arrays. \\
Libraries/OneWire & Reset, bit slots, bytes, Dallas CRC8 and ROM search. \\
Libraries/DS18B20 & Addressed commands, scratchpad decoding and array scheduler. \\
Drivers & Local vendor HAL and CMSIS shared by both build systems\cite{cubef1}. \\
MDK-ARM & Keil project and Arm assembly startup. \\
Core/Startup & GNU assembly startup for PlatformIO. \\
tools and tests & Regeneration synchronization, checks and emulated C tests. \\
\bottomrule
\end{tabular}

\subsection{The hardware contract}
\texttt{OneWire\_Bus} binds a context pointer and callbacks. Drive-low asserts
the line; release makes it high impedance; read samples the pin. Delay accepts
microseconds and must run while interrupts are disabled. The optional critical
callbacks must be supplied as a pair: entering returns previous interrupt state,
leaving restores that exact state. Supplying neither requires the application's
environment to guarantee bit timing independently.

On this Cortex-M3, DWT cycle counting supplies the delay. The adapter enables
trace and cycle counting after the system clock is configured. GPIO BSRR/BRR
registers perform single writes without GPIO reconfiguration. The elapsed-cycle
subtraction handles counter wrap for these short delays. Other MCUs, including
ARM families without a usable DWT counter, should implement a timer callback.
The portable headers include only standard C headers and other portable headers.

\subsection{Ownership and concurrency}
The caller allocates the array object, ROM array, discovery buffer, temperatures
and per-slot status. Both ROM buffers have the configured capacity, are separate,
and remain alive throughout use. No allocation or global library singleton is
required. Different buses may have separate instances. Calls on the same bus
must be serialized by the application; never invoke the driver from an ISR or
change core frequency during a transaction. An RTOS mutex belongs around complete
transactions, not around an individual callback.
